\section{区間の両端を交換する}\label{節:区間の反転}

対称の中心は原点に限らない。区間 $[a,b]$ の両端を交換するには，$x=a+b-t$ と置けばよい。対応は
\[
x=a\longleftrightarrow t=b,\qquad
x=b\longleftrightarrow t=a
\]
であり，$\dd x=-\dd t$ だから，
\begin{equation}\label{式:区間反転}
\int_a^b f(x)\dd x=\int_a^b f(a+b-x)\dd x
\end{equation}
を得る。置換の直後は積分の上下端が逆になるが，微分の負号と合わせて向きを戻している。

この式の左右は同じ値なので，加えて二で割ると
\begin{equation}\label{式:反転平均}
\int_a^b f(x)\dd x
 =\frac12\int_a^b\{f(x)+f(a+b-x)\}\dd x
\end{equation}
となる。$f(x)+f(a+b-x)$ が定数 $C$ になれば，答えは $C(b-a)$ の半分である。元の式を個別に積分できるかどうかは，この計算では問題にならない。

今度は，関数 $g$ が区間の中央について対称，すなわち
\[
g(a+b-x)=g(x)
\]
を満たすとする。このとき $xg(x)$ の積分を $I$ とすると，反転により
\[
I=\int_a^b(a+b-x)g(x)\dd x.
\]
元の式を加えると $x+(a+b-x)=a+b$ なので，
\begin{equation}\label{式:対称な重みの一次因子}
\int_a^b xg(x)\dd x
 =\frac{a+b}{2}\int_a^b g(x)\dd x
\end{equation}
となる。中央からの距離が等しい二点では $g$ の値が等しく，その二点の $x$ 座標の平均は常に区間の中央になる，という関係である。

\begin{example}{変数が掛かった三角関数}
次の定積分を求めよ。
\[
I=\int_0^\pi x\sin^3x\dd x.
\]
\end{example}

$x$ と三角関数の積を見ると部分積分を考えたくなる。しかしこの区間では $\sin(\pi-x)=\sin x$ である。$g(x)=\sin^3x$ は区間の中央について対称なので，式\eqref{式:対称な重みの一次因子}から
\[
I=\frac{\pi}{2}\int_0^\pi\sin^3x\dd x.
\]
残りは $\sin^3x=(1-\cos^2x)\sin x$ として $u=\cos x$ と置く。
\[
\int_0^\pi\sin^3x\dd x
 =\int_{-1}^1(1-u^2)\dd u
 =\left[u-\frac{u^3}{3}\right]_{-1}^1
 =\frac43.
\]
したがって $I=\frac{2\pi}{3}$ である。反転で $x$ を取り除き，置換で三角関数を取り除くという二段階で計算できた。

式\eqref{式:対称な重みの一次因子}を使う際には，$xg(x)$ 全体が対称なのではなく，$g$ が対称であることを確認する。たとえば $\int_0^1 x\e^x\dd x$ では $\e^{1-x}=\e^x$ が成り立たないので，$x$ を区間の中央 $\frac12$ に置き換えることはできない。

また，$x^2g(x)$ まで同じように中央の二乗に置き換えられるわけでもない。$x^2+(a+b-x)^2$ は $x$ を含むからである。反転して足すことで何が消えるのかを，その都度式で確かめる必要がある。

